% ============================================================
%   Chapter 8  \textemdash\  Learning Collision Avoidance from
%                            LLM-Derived Rewards
% ------------------------------------------------------------
%   DRAFT SCAFFOLD (design-first). Not yet wired into main.tex.
%   To include: uncomment the % ============================================================
%   Chapter 8  \textemdash\  Learning Collision Avoidance from
%                            LLM-Derived Rewards
% ------------------------------------------------------------
%   DRAFT SCAFFOLD (design-first). Not yet wired into main.tex.
%   To include: uncomment the % ============================================================
%   Chapter 8  \textemdash\  Learning Collision Avoidance from
%                            LLM-Derived Rewards
% ------------------------------------------------------------
%   DRAFT SCAFFOLD (design-first). Not yet wired into main.tex.
%   To include: uncomment the % ============================================================
%   Chapter 8  \textemdash\  Learning Collision Avoidance from
%                            LLM-Derived Rewards
% ------------------------------------------------------------
%   DRAFT SCAFFOLD (design-first). Not yet wired into main.tex.
%   To include: uncomment the \input{tex/08_llm_reward_exposure}
%   line in main.tex AFTER the figures listed below exist in
%   Memoria/imgs/ and the bib keys in §10 of
%   docs/shielded_motif_design.md are added to referencias.bib.
%
%   Planned figures (imgs/, generate from run SQLite as in Ch.7):
%     - encounter_rate_curve   (exposure director lifts encounter_rate ~3% -> ~30%)
%     - reward_field           (r_phi over front-range x closing-speed slice)
%     - arms_crash_bars        (C0/C1/C2/C3 shield-OFF crash @60 NPC)
%     - shaping_learning_curve (shield-OFF crash vs updates, per arm)
%   Planned tables:
%     - tab:cov-arms           (factorial arms x outcomes)
% ============================================================

\chapter{Learning Collision Avoidance from LLM-Derived Rewards}\label{ch:llmreward}

Chapter~\ref{ch:dependence} closed on an explicit open problem. Behavioural cloning of the
shield's steering removed shield \emph{dependence} \textemdash\ the weaned agent drives as well
unaided as it does shielded \textemdash\ but it left a residual \emph{collision} gap: without the
shield the agent still crashes on roughly one episode in six, because the shield's throttle is an
emergency brake and is therefore a poor longitudinal teacher (Section~\ref{sec:dep-collapse}).
That chapter argued the gap is ``better posed as a reward-shaping problem \dots\ that lets the
agent learn its \emph{own} gentle avoidance''. This chapter takes up that problem, and in doing so
revisits \textemdash\ and finally vindicates \textemdash\ the use of a Large Language Model in the
training loop, after the negative result of Section~\ref{sec:dep-llm}.

The chapter is organised around a single thesis: the earlier LLM failure and the failure of a
hand-designed time-to-collision (TTC) penalty share \emph{one} root cause, and that cause dictates
the design that finally works. Section~\ref{sec:cov-diagnosis} establishes the diagnosis;
Section~\ref{sec:cov-method} derives the method (an LLM-learned dense reward trained on an
LLM-generated exposure curriculum); Sections~\ref{sec:cov-exposure}--\ref{sec:cov-shaping} detail
its three components; Section~\ref{sec:cov-results} reports the experiments; and
Section~\ref{sec:cov-discussion} places the result in the \emph{experience\,>\,gradient\,>\,
incentive} hierarchy that organises the whole of Chapters~\ref{ch:dependence}--\ref{ch:llmreward}.

\section{Diagnosis: one root cause behind two negative results}\label{sec:cov-diagnosis}

Two attempts to use a language model or a designed reward to close the avoidance gap failed, and
they failed identically.

\paragraph{The permissiveness dial (Section~\ref{sec:dep-llm}).} The LLM controlled a single
scalar every hundred episodes; its output was \emph{not causally coupled} to the metric it was
meant to move (the correlation between permissiveness and shield-activation rate was
\(\approx 0\)). A fixed schedule would have done the same.

\paragraph{The TTC penalty.} A hand-designed time-to-collision / following-distance penalty was
added to the reward shaper and its weight swept across four runs. The penalty \emph{fired in only
\(\approx 3\%\) of episodes}, its magnitude when fired was flat across weights, and the shield-off
crash rate did not fall \textemdash\ it merely redistributed failures from off-road to collision.
% TODO: cite/forward-ref the TTC sweep table if it is written up elsewhere.

\paragraph{The common cause.} Both signals were starved. On the Town04 highway loop, globally
spawned traffic places a vehicle in the ego's front cone in only about \(3\%\) of episodes, so any
collision-specific signal \textemdash\ whether an LLM dial or a TTC penalty \textemdash\ almost
never has anything to act on. Three constraints follow, and the method of
Section~\ref{sec:cov-method} is the minimal design that satisfies all three:
\begin{enumerate}
  \item \textbf{Bandwidth.} The learning signal must be \emph{dense} (per step), like the
        behavioural-cloning term of Equation~\eqref{eq:bc}, not a per-window scalar.
  \item \textbf{Exposure.} A collision signal starves at \(3\%\) coverage; the relevant
        experience must be \emph{manufactured}.
  \item \textbf{Stationarity.} A live additive reward rescales the running return
        normaliser the critic is calibrated against (Section~\ref{sec:design-agent-ppo}); the
        signal must be \emph{policy-invariant}.
\end{enumerate}

\section{Method: dense LLM reward over generated exposure}\label{sec:cov-method}

The proposed method, which we call \emph{Shielded Motif under generated exposure}, has three
components that map one-to-one onto the three constraints. An \textbf{exposure director}
(Section~\ref{sec:cov-exposure}) manufactures vehicle encounters, supplying the missing
experience. A \textbf{preference-trained reward model} (Section~\ref{sec:cov-reward}) supplies a
dense, per-step safety signal learned from a language model's judgements \textemdash\ offline,
so the model never runs inside the control loop. A \textbf{potential-based shaping}
integration (Section~\ref{sec:cov-shaping}) delivers that signal without altering the optimal
policy. Throughout, the safety shield and the deterministic safety governor of
Chapter~\ref{ch:dependence} are retained unchanged: the shield still prevents unsafe
\emph{exploration} on the harder, manufactured scenarios, and the ``proposer~+~deterministic
veto'' pattern \textemdash\ already shown to be sound independently of the LLM
(Section~\ref{sec:dep-discussion}) \textemdash\ now guards the exposure curriculum.

The design follows the Motif paradigm of intrinsic motivation from artificial-intelligence
feedback~\cite{klissarov2023motif}, in which an LLM expresses preferences over pairs of
observation captions that are distilled into a reward model; it inherits the preference-learning
formulation of~\cite{christiano2017preferences} and the policy-invariance guarantee
of~\cite{ng1999shaping}. It is related to, but distinct from, language-conditioned reward
generation for driving~\cite{huang2025vlmrl} and to LLM-based safety-critical scenario
generation~\cite{zhang2024chatscene,nie2025seeking}, from which the exposure director borrows its
intent-conditioned design.

\section{Component I \textemdash\ the exposure director}\label{sec:cov-exposure}

% Narrative to write:
%  - route-traffic injection along the ego waypoint chain (leads + cut-ins),
%    absolute desired speeds (not relative to the speed limit), governor-gated.
%  - failure-driven scenario generation: shield-off crash traces -> LLM infers a
%    behavioural intent -> structured scenario spec; intent->recipe memory bank
%    (after nie2025seeking, zhang2024chatscene).
%  - progress KPI: encounter_rate = fraction of episodes with min_front_dynamic < 0.7.
%  - safety: a dedicated governor with exposure ceilings (crash 0.20 / target 0.12)
%    that CUTS traffic on a crash breach.
%  - figure: encounter_rate_curve.

\section{Component II \textemdash\ the preference-trained reward model}\label{sec:cov-reward}

% Narrative to write:
%  - offline pipeline: sample observations (focus on the dynamic-LIDAR channel
%    obs[240:480] + front-cone scalars), template each into a structured caption,
%    LLM ranks PAIRS into preferences (reuse the schema-constrained, deterministic,
%    fallback-guarded decoding of the abandoned shield tuner).
%  - reward model r_phi: MLP 240(+k)->64->64->1; Bradley-Terry loss below.
%  - DISTINCTION from a crash-label safety critic: dense graded preference signal
%    available on near-miss steps that never crashed, vs sparse binary crash labels.

The reward model \(r_\phi\) is trained from preferences with the Bradley--Terry
objective~\cite{christiano2017preferences}: for an annotated pair of scenes \((s_a, s_b)\) in
which the language model judges \(s_a\) the safer situation,
\begin{equation}\label{eq:bt}
  \mathcal{L}(\phi) \;=\; -\,\mathbb{E}_{(s_a,s_b)}\;
  c_{ab}\,\log \sigma\!\bigl(r_\phi(s_a) - r_\phi(s_b)\bigr),
\end{equation}
where \(\sigma\) is the logistic function and \(c_{ab}\in[0,1]\) is the language model's stated
confidence, used to down-weight uncertain pairs. The annotation is performed entirely offline, so
the language model never contends with CARLA for GPU memory and adds no latency to the \(20\,
\mathrm{Hz}\) control loop \textemdash\ directly removing the operational difficulty that the
in-loop controller of Section~\ref{sec:dep-llm} suffered.

\section{Component III \textemdash\ policy-invariant shaping}\label{sec:cov-shaping}

The learned reward is delivered as potential-based shaping~\cite{ng1999shaping}, using \(r_\phi\)
as the potential:
\begin{equation}\label{eq:potential}
  F_t \;=\; \gamma\,r_\phi(s_{t+1}) - r_\phi(s_t),
  \qquad
  r^{\mathrm{shaped}}_t \;=\; r_t \;+\; \beta\,F_t .
\end{equation}
Because \(F_t\) is a difference of potentials, the set of optimal policies is provably
unchanged~\cite{ng1999shaping}: the term can only \emph{re-time} credit toward the states that
precede danger, it cannot create a new optimum. This is the formal guarantee that the
park-attractor of Section~\ref{sec:dep-collapse} \textemdash\ in which a direct supervised
gradient overpowered the reward and drove the policy to a stationary equilibrium \textemdash\
cannot recur here, and it is the reason the design respects the reward-stationarity constraint
that the original meta-controller was deliberately built to honour. The coefficient \(\beta\) is
frozen per run and annealed toward a small floor, strong early to install the avoidance skill.

\section{Experiments}\label{sec:cov-results}

% Pre-registered factorial design (see docs/shielded_motif_design.md, table tab:cov-arms):
%  C0 control (= learn_fix2), C1 exposure-only, C2 reward-only (predicted to STARVE),
%  C3 full. All warm-start learn_fix2 with --log_std_anchor_coef 1e-3 (mandatory).
%  Primary metric: shield-OFF crash @60 NPC (main_eval --num_npc 60 --seed 100).
%  Success: C3 shield-OFF crash < 16% (the learn_fix2 control).
%  Decisive/falsification test: if exposure lifts encounter_rate to ~30% and crash
%  still will not move, that is a clean, causally-isolated third negative result.

\begin{table}[H]
  \centering
  \caption{Pre-registered factorial arms. Each isolates one leg of the method; C2 is predicted
           to \emph{starve} (a reward with no exposure), which would itself confirm the
           exposure diagnosis of Section~\ref{sec:cov-diagnosis}.}
  \label{tab:cov-arms}
  \small
  \begin{tabular}{@{}llll@{}}
    \toprule
    \textbf{Arm} & \textbf{Exposure} & \textbf{Pref.\ reward} & \textbf{Purpose} \\
    \midrule
    C0 control       & global NPC only    & off & reproduce \texttt{learn\_fix2} \\
    C1 exposure-only & LLM director       & off & does exposure alone move crash? \\
    C2 reward-only   & global NPC only    & on  & predicted to starve (3\% coverage) \\
    C3 full          & LLM director       & on  & the proposed method \\
    \bottomrule
  \end{tabular}
\end{table}

% TODO figures once runs land:
% \begin{figure}[H]\centering\includegraphics[width=0.85\textwidth]{imgs/arms_crash_bars}
%   \caption{Shield-off crash rate at 60 NPCs across the four arms.}\label{fig:cov-arms}\end{figure}

\section{Discussion}\label{sec:cov-discussion}

% Narrative to write:
%  - the experience > gradient > incentive hierarchy: TTC (incentive) failed for lack of
%    exposure; shield BC (gradient) fixed steering but not throttle; only when an LLM
%    manufactures the experience does a learned dense reward teach self-driven avoidance.
%  - salvage story: the abandoned meta-stack (schema-constrained decoding, deterministic
%    governor, async runner) is reused as trustworthy offline annotation infrastructure;
%    the negative result of Ch.7 is what FORCED this (correct) design, not a lucky guess.
%  - honest limitation: if the decisive test fails, report the bounded negative result.

\clearpage

%   line in main.tex AFTER the figures listed below exist in
%   Memoria/imgs/ and the bib keys in §10 of
%   docs/shielded_motif_design.md are added to referencias.bib.
%
%   Planned figures (imgs/, generate from run SQLite as in Ch.7):
%     - encounter_rate_curve   (exposure director lifts encounter_rate ~3% -> ~30%)
%     - reward_field           (r_phi over front-range x closing-speed slice)
%     - arms_crash_bars        (C0/C1/C2/C3 shield-OFF crash @60 NPC)
%     - shaping_learning_curve (shield-OFF crash vs updates, per arm)
%   Planned tables:
%     - tab:cov-arms           (factorial arms x outcomes)
% ============================================================

\chapter{Learning Collision Avoidance from LLM-Derived Rewards}\label{ch:llmreward}

Chapter~\ref{ch:dependence} closed on an explicit open problem. Behavioural cloning of the
shield's steering removed shield \emph{dependence} \textemdash\ the weaned agent drives as well
unaided as it does shielded \textemdash\ but it left a residual \emph{collision} gap: without the
shield the agent still crashes on roughly one episode in six, because the shield's throttle is an
emergency brake and is therefore a poor longitudinal teacher (Section~\ref{sec:dep-collapse}).
That chapter argued the gap is ``better posed as a reward-shaping problem \dots\ that lets the
agent learn its \emph{own} gentle avoidance''. This chapter takes up that problem, and in doing so
revisits \textemdash\ and finally vindicates \textemdash\ the use of a Large Language Model in the
training loop, after the negative result of Section~\ref{sec:dep-llm}.

The chapter is organised around a single thesis: the earlier LLM failure and the failure of a
hand-designed time-to-collision (TTC) penalty share \emph{one} root cause, and that cause dictates
the design that finally works. Section~\ref{sec:cov-diagnosis} establishes the diagnosis;
Section~\ref{sec:cov-method} derives the method (an LLM-learned dense reward trained on an
LLM-generated exposure curriculum); Sections~\ref{sec:cov-exposure}--\ref{sec:cov-shaping} detail
its three components; Section~\ref{sec:cov-results} reports the experiments; and
Section~\ref{sec:cov-discussion} places the result in the \emph{experience\,>\,gradient\,>\,
incentive} hierarchy that organises the whole of Chapters~\ref{ch:dependence}--\ref{ch:llmreward}.

\section{Diagnosis: one root cause behind two negative results}\label{sec:cov-diagnosis}

Two attempts to use a language model or a designed reward to close the avoidance gap failed, and
they failed identically.

\paragraph{The permissiveness dial (Section~\ref{sec:dep-llm}).} The LLM controlled a single
scalar every hundred episodes; its output was \emph{not causally coupled} to the metric it was
meant to move (the correlation between permissiveness and shield-activation rate was
\(\approx 0\)). A fixed schedule would have done the same.

\paragraph{The TTC penalty.} A hand-designed time-to-collision / following-distance penalty was
added to the reward shaper and its weight swept across four runs. The penalty \emph{fired in only
\(\approx 3\%\) of episodes}, its magnitude when fired was flat across weights, and the shield-off
crash rate did not fall \textemdash\ it merely redistributed failures from off-road to collision.
% TODO: cite/forward-ref the TTC sweep table if it is written up elsewhere.

\paragraph{The common cause.} Both signals were starved. On the Town04 highway loop, globally
spawned traffic places a vehicle in the ego's front cone in only about \(3\%\) of episodes, so any
collision-specific signal \textemdash\ whether an LLM dial or a TTC penalty \textemdash\ almost
never has anything to act on. Three constraints follow, and the method of
Section~\ref{sec:cov-method} is the minimal design that satisfies all three:
\begin{enumerate}
  \item \textbf{Bandwidth.} The learning signal must be \emph{dense} (per step), like the
        behavioural-cloning term of Equation~\eqref{eq:bc}, not a per-window scalar.
  \item \textbf{Exposure.} A collision signal starves at \(3\%\) coverage; the relevant
        experience must be \emph{manufactured}.
  \item \textbf{Stationarity.} A live additive reward rescales the running return
        normaliser the critic is calibrated against (Section~\ref{sec:design-agent-ppo}); the
        signal must be \emph{policy-invariant}.
\end{enumerate}

\section{Method: dense LLM reward over generated exposure}\label{sec:cov-method}

The proposed method, which we call \emph{Shielded Motif under generated exposure}, has three
components that map one-to-one onto the three constraints. An \textbf{exposure director}
(Section~\ref{sec:cov-exposure}) manufactures vehicle encounters, supplying the missing
experience. A \textbf{preference-trained reward model} (Section~\ref{sec:cov-reward}) supplies a
dense, per-step safety signal learned from a language model's judgements \textemdash\ offline,
so the model never runs inside the control loop. A \textbf{potential-based shaping}
integration (Section~\ref{sec:cov-shaping}) delivers that signal without altering the optimal
policy. Throughout, the safety shield and the deterministic safety governor of
Chapter~\ref{ch:dependence} are retained unchanged: the shield still prevents unsafe
\emph{exploration} on the harder, manufactured scenarios, and the ``proposer~+~deterministic
veto'' pattern \textemdash\ already shown to be sound independently of the LLM
(Section~\ref{sec:dep-discussion}) \textemdash\ now guards the exposure curriculum.

The design follows the Motif paradigm of intrinsic motivation from artificial-intelligence
feedback~\cite{klissarov2023motif}, in which an LLM expresses preferences over pairs of
observation captions that are distilled into a reward model; it inherits the preference-learning
formulation of~\cite{christiano2017preferences} and the policy-invariance guarantee
of~\cite{ng1999shaping}. It is related to, but distinct from, language-conditioned reward
generation for driving~\cite{huang2025vlmrl} and to LLM-based safety-critical scenario
generation~\cite{zhang2024chatscene,nie2025seeking}, from which the exposure director borrows its
intent-conditioned design.

\section{Component I \textemdash\ the exposure director}\label{sec:cov-exposure}

% Narrative to write:
%  - route-traffic injection along the ego waypoint chain (leads + cut-ins),
%    absolute desired speeds (not relative to the speed limit), governor-gated.
%  - failure-driven scenario generation: shield-off crash traces -> LLM infers a
%    behavioural intent -> structured scenario spec; intent->recipe memory bank
%    (after nie2025seeking, zhang2024chatscene).
%  - progress KPI: encounter_rate = fraction of episodes with min_front_dynamic < 0.7.
%  - safety: a dedicated governor with exposure ceilings (crash 0.20 / target 0.12)
%    that CUTS traffic on a crash breach.
%  - figure: encounter_rate_curve.

\section{Component II \textemdash\ the preference-trained reward model}\label{sec:cov-reward}

% Narrative to write:
%  - offline pipeline: sample observations (focus on the dynamic-LIDAR channel
%    obs[240:480] + front-cone scalars), template each into a structured caption,
%    LLM ranks PAIRS into preferences (reuse the schema-constrained, deterministic,
%    fallback-guarded decoding of the abandoned shield tuner).
%  - reward model r_phi: MLP 240(+k)->64->64->1; Bradley-Terry loss below.
%  - DISTINCTION from a crash-label safety critic: dense graded preference signal
%    available on near-miss steps that never crashed, vs sparse binary crash labels.

The reward model \(r_\phi\) is trained from preferences with the Bradley--Terry
objective~\cite{christiano2017preferences}: for an annotated pair of scenes \((s_a, s_b)\) in
which the language model judges \(s_a\) the safer situation,
\begin{equation}\label{eq:bt}
  \mathcal{L}(\phi) \;=\; -\,\mathbb{E}_{(s_a,s_b)}\;
  c_{ab}\,\log \sigma\!\bigl(r_\phi(s_a) - r_\phi(s_b)\bigr),
\end{equation}
where \(\sigma\) is the logistic function and \(c_{ab}\in[0,1]\) is the language model's stated
confidence, used to down-weight uncertain pairs. The annotation is performed entirely offline, so
the language model never contends with CARLA for GPU memory and adds no latency to the \(20\,
\mathrm{Hz}\) control loop \textemdash\ directly removing the operational difficulty that the
in-loop controller of Section~\ref{sec:dep-llm} suffered.

\section{Component III \textemdash\ policy-invariant shaping}\label{sec:cov-shaping}

The learned reward is delivered as potential-based shaping~\cite{ng1999shaping}, using \(r_\phi\)
as the potential:
\begin{equation}\label{eq:potential}
  F_t \;=\; \gamma\,r_\phi(s_{t+1}) - r_\phi(s_t),
  \qquad
  r^{\mathrm{shaped}}_t \;=\; r_t \;+\; \beta\,F_t .
\end{equation}
Because \(F_t\) is a difference of potentials, the set of optimal policies is provably
unchanged~\cite{ng1999shaping}: the term can only \emph{re-time} credit toward the states that
precede danger, it cannot create a new optimum. This is the formal guarantee that the
park-attractor of Section~\ref{sec:dep-collapse} \textemdash\ in which a direct supervised
gradient overpowered the reward and drove the policy to a stationary equilibrium \textemdash\
cannot recur here, and it is the reason the design respects the reward-stationarity constraint
that the original meta-controller was deliberately built to honour. The coefficient \(\beta\) is
frozen per run and annealed toward a small floor, strong early to install the avoidance skill.

\section{Experiments}\label{sec:cov-results}

% Pre-registered factorial design (see docs/shielded_motif_design.md, table tab:cov-arms):
%  C0 control (= learn_fix2), C1 exposure-only, C2 reward-only (predicted to STARVE),
%  C3 full. All warm-start learn_fix2 with --log_std_anchor_coef 1e-3 (mandatory).
%  Primary metric: shield-OFF crash @60 NPC (main_eval --num_npc 60 --seed 100).
%  Success: C3 shield-OFF crash < 16% (the learn_fix2 control).
%  Decisive/falsification test: if exposure lifts encounter_rate to ~30% and crash
%  still will not move, that is a clean, causally-isolated third negative result.

\begin{table}[H]
  \centering
  \caption{Pre-registered factorial arms. Each isolates one leg of the method; C2 is predicted
           to \emph{starve} (a reward with no exposure), which would itself confirm the
           exposure diagnosis of Section~\ref{sec:cov-diagnosis}.}
  \label{tab:cov-arms}
  \small
  \begin{tabular}{@{}llll@{}}
    \toprule
    \textbf{Arm} & \textbf{Exposure} & \textbf{Pref.\ reward} & \textbf{Purpose} \\
    \midrule
    C0 control       & global NPC only    & off & reproduce \texttt{learn\_fix2} \\
    C1 exposure-only & LLM director       & off & does exposure alone move crash? \\
    C2 reward-only   & global NPC only    & on  & predicted to starve (3\% coverage) \\
    C3 full          & LLM director       & on  & the proposed method \\
    \bottomrule
  \end{tabular}
\end{table}

% TODO figures once runs land:
% \begin{figure}[H]\centering\includegraphics[width=0.85\textwidth]{imgs/arms_crash_bars}
%   \caption{Shield-off crash rate at 60 NPCs across the four arms.}\label{fig:cov-arms}\end{figure}

\section{Discussion}\label{sec:cov-discussion}

% Narrative to write:
%  - the experience > gradient > incentive hierarchy: TTC (incentive) failed for lack of
%    exposure; shield BC (gradient) fixed steering but not throttle; only when an LLM
%    manufactures the experience does a learned dense reward teach self-driven avoidance.
%  - salvage story: the abandoned meta-stack (schema-constrained decoding, deterministic
%    governor, async runner) is reused as trustworthy offline annotation infrastructure;
%    the negative result of Ch.7 is what FORCED this (correct) design, not a lucky guess.
%  - honest limitation: if the decisive test fails, report the bounded negative result.

\clearpage

%   line in main.tex AFTER the figures listed below exist in
%   Memoria/imgs/ and the bib keys in §10 of
%   docs/shielded_motif_design.md are added to referencias.bib.
%
%   Planned figures (imgs/, generate from run SQLite as in Ch.7):
%     - encounter_rate_curve   (exposure director lifts encounter_rate ~3% -> ~30%)
%     - reward_field           (r_phi over front-range x closing-speed slice)
%     - arms_crash_bars        (C0/C1/C2/C3 shield-OFF crash @60 NPC)
%     - shaping_learning_curve (shield-OFF crash vs updates, per arm)
%   Planned tables:
%     - tab:cov-arms           (factorial arms x outcomes)
% ============================================================

\chapter{Learning Collision Avoidance from LLM-Derived Rewards}\label{ch:llmreward}

Chapter~\ref{ch:dependence} closed on an explicit open problem. Behavioural cloning of the
shield's steering removed shield \emph{dependence} \textemdash\ the weaned agent drives as well
unaided as it does shielded \textemdash\ but it left a residual \emph{collision} gap: without the
shield the agent still crashes on roughly one episode in six, because the shield's throttle is an
emergency brake and is therefore a poor longitudinal teacher (Section~\ref{sec:dep-collapse}).
That chapter argued the gap is ``better posed as a reward-shaping problem \dots\ that lets the
agent learn its \emph{own} gentle avoidance''. This chapter takes up that problem, and in doing so
revisits \textemdash\ and finally vindicates \textemdash\ the use of a Large Language Model in the
training loop, after the negative result of Section~\ref{sec:dep-llm}.

The chapter is organised around a single thesis: the earlier LLM failure and the failure of a
hand-designed time-to-collision (TTC) penalty share \emph{one} root cause, and that cause dictates
the design that finally works. Section~\ref{sec:cov-diagnosis} establishes the diagnosis;
Section~\ref{sec:cov-method} derives the method (an LLM-learned dense reward trained on an
LLM-generated exposure curriculum); Sections~\ref{sec:cov-exposure}--\ref{sec:cov-shaping} detail
its three components; Section~\ref{sec:cov-results} reports the experiments; and
Section~\ref{sec:cov-discussion} places the result in the \emph{experience\,>\,gradient\,>\,
incentive} hierarchy that organises the whole of Chapters~\ref{ch:dependence}--\ref{ch:llmreward}.

\section{Diagnosis: one root cause behind two negative results}\label{sec:cov-diagnosis}

Two attempts to use a language model or a designed reward to close the avoidance gap failed, and
they failed identically.

\paragraph{The permissiveness dial (Section~\ref{sec:dep-llm}).} The LLM controlled a single
scalar every hundred episodes; its output was \emph{not causally coupled} to the metric it was
meant to move (the correlation between permissiveness and shield-activation rate was
\(\approx 0\)). A fixed schedule would have done the same.

\paragraph{The TTC penalty.} A hand-designed time-to-collision / following-distance penalty was
added to the reward shaper and its weight swept across four runs. The penalty \emph{fired in only
\(\approx 3\%\) of episodes}, its magnitude when fired was flat across weights, and the shield-off
crash rate did not fall \textemdash\ it merely redistributed failures from off-road to collision.
% TODO: cite/forward-ref the TTC sweep table if it is written up elsewhere.

\paragraph{The common cause.} Both signals were starved. On the Town04 highway loop, globally
spawned traffic places a vehicle in the ego's front cone in only about \(3\%\) of episodes, so any
collision-specific signal \textemdash\ whether an LLM dial or a TTC penalty \textemdash\ almost
never has anything to act on. Three constraints follow, and the method of
Section~\ref{sec:cov-method} is the minimal design that satisfies all three:
\begin{enumerate}
  \item \textbf{Bandwidth.} The learning signal must be \emph{dense} (per step), like the
        behavioural-cloning term of Equation~\eqref{eq:bc}, not a per-window scalar.
  \item \textbf{Exposure.} A collision signal starves at \(3\%\) coverage; the relevant
        experience must be \emph{manufactured}.
  \item \textbf{Stationarity.} A live additive reward rescales the running return
        normaliser the critic is calibrated against (Section~\ref{sec:design-agent-ppo}); the
        signal must be \emph{policy-invariant}.
\end{enumerate}

\section{Method: dense LLM reward over generated exposure}\label{sec:cov-method}

The proposed method, which we call \emph{Shielded Motif under generated exposure}, has three
components that map one-to-one onto the three constraints. An \textbf{exposure director}
(Section~\ref{sec:cov-exposure}) manufactures vehicle encounters, supplying the missing
experience. A \textbf{preference-trained reward model} (Section~\ref{sec:cov-reward}) supplies a
dense, per-step safety signal learned from a language model's judgements \textemdash\ offline,
so the model never runs inside the control loop. A \textbf{potential-based shaping}
integration (Section~\ref{sec:cov-shaping}) delivers that signal without altering the optimal
policy. Throughout, the safety shield and the deterministic safety governor of
Chapter~\ref{ch:dependence} are retained unchanged: the shield still prevents unsafe
\emph{exploration} on the harder, manufactured scenarios, and the ``proposer~+~deterministic
veto'' pattern \textemdash\ already shown to be sound independently of the LLM
(Section~\ref{sec:dep-discussion}) \textemdash\ now guards the exposure curriculum.

The design follows the Motif paradigm of intrinsic motivation from artificial-intelligence
feedback~\cite{klissarov2023motif}, in which an LLM expresses preferences over pairs of
observation captions that are distilled into a reward model; it inherits the preference-learning
formulation of~\cite{christiano2017preferences} and the policy-invariance guarantee
of~\cite{ng1999shaping}. It is related to, but distinct from, language-conditioned reward
generation for driving~\cite{huang2025vlmrl} and to LLM-based safety-critical scenario
generation~\cite{zhang2024chatscene,nie2025seeking}, from which the exposure director borrows its
intent-conditioned design.

\section{Component I \textemdash\ the exposure director}\label{sec:cov-exposure}

% Narrative to write:
%  - route-traffic injection along the ego waypoint chain (leads + cut-ins),
%    absolute desired speeds (not relative to the speed limit), governor-gated.
%  - failure-driven scenario generation: shield-off crash traces -> LLM infers a
%    behavioural intent -> structured scenario spec; intent->recipe memory bank
%    (after nie2025seeking, zhang2024chatscene).
%  - progress KPI: encounter_rate = fraction of episodes with min_front_dynamic < 0.7.
%  - safety: a dedicated governor with exposure ceilings (crash 0.20 / target 0.12)
%    that CUTS traffic on a crash breach.
%  - figure: encounter_rate_curve.

\section{Component II \textemdash\ the preference-trained reward model}\label{sec:cov-reward}

% Narrative to write:
%  - offline pipeline: sample observations (focus on the dynamic-LIDAR channel
%    obs[240:480] + front-cone scalars), template each into a structured caption,
%    LLM ranks PAIRS into preferences (reuse the schema-constrained, deterministic,
%    fallback-guarded decoding of the abandoned shield tuner).
%  - reward model r_phi: MLP 240(+k)->64->64->1; Bradley-Terry loss below.
%  - DISTINCTION from a crash-label safety critic: dense graded preference signal
%    available on near-miss steps that never crashed, vs sparse binary crash labels.

The reward model \(r_\phi\) is trained from preferences with the Bradley--Terry
objective~\cite{christiano2017preferences}: for an annotated pair of scenes \((s_a, s_b)\) in
which the language model judges \(s_a\) the safer situation,
\begin{equation}\label{eq:bt}
  \mathcal{L}(\phi) \;=\; -\,\mathbb{E}_{(s_a,s_b)}\;
  c_{ab}\,\log \sigma\!\bigl(r_\phi(s_a) - r_\phi(s_b)\bigr),
\end{equation}
where \(\sigma\) is the logistic function and \(c_{ab}\in[0,1]\) is the language model's stated
confidence, used to down-weight uncertain pairs. The annotation is performed entirely offline, so
the language model never contends with CARLA for GPU memory and adds no latency to the \(20\,
\mathrm{Hz}\) control loop \textemdash\ directly removing the operational difficulty that the
in-loop controller of Section~\ref{sec:dep-llm} suffered.

\section{Component III \textemdash\ policy-invariant shaping}\label{sec:cov-shaping}

The learned reward is delivered as potential-based shaping~\cite{ng1999shaping}, using \(r_\phi\)
as the potential:
\begin{equation}\label{eq:potential}
  F_t \;=\; \gamma\,r_\phi(s_{t+1}) - r_\phi(s_t),
  \qquad
  r^{\mathrm{shaped}}_t \;=\; r_t \;+\; \beta\,F_t .
\end{equation}
Because \(F_t\) is a difference of potentials, the set of optimal policies is provably
unchanged~\cite{ng1999shaping}: the term can only \emph{re-time} credit toward the states that
precede danger, it cannot create a new optimum. This is the formal guarantee that the
park-attractor of Section~\ref{sec:dep-collapse} \textemdash\ in which a direct supervised
gradient overpowered the reward and drove the policy to a stationary equilibrium \textemdash\
cannot recur here, and it is the reason the design respects the reward-stationarity constraint
that the original meta-controller was deliberately built to honour. The coefficient \(\beta\) is
frozen per run and annealed toward a small floor, strong early to install the avoidance skill.

\section{Experiments}\label{sec:cov-results}

% Pre-registered factorial design (see docs/shielded_motif_design.md, table tab:cov-arms):
%  C0 control (= learn_fix2), C1 exposure-only, C2 reward-only (predicted to STARVE),
%  C3 full. All warm-start learn_fix2 with --log_std_anchor_coef 1e-3 (mandatory).
%  Primary metric: shield-OFF crash @60 NPC (main_eval --num_npc 60 --seed 100).
%  Success: C3 shield-OFF crash < 16% (the learn_fix2 control).
%  Decisive/falsification test: if exposure lifts encounter_rate to ~30% and crash
%  still will not move, that is a clean, causally-isolated third negative result.

\begin{table}[H]
  \centering
  \caption{Pre-registered factorial arms. Each isolates one leg of the method; C2 is predicted
           to \emph{starve} (a reward with no exposure), which would itself confirm the
           exposure diagnosis of Section~\ref{sec:cov-diagnosis}.}
  \label{tab:cov-arms}
  \small
  \begin{tabular}{@{}llll@{}}
    \toprule
    \textbf{Arm} & \textbf{Exposure} & \textbf{Pref.\ reward} & \textbf{Purpose} \\
    \midrule
    C0 control       & global NPC only    & off & reproduce \texttt{learn\_fix2} \\
    C1 exposure-only & LLM director       & off & does exposure alone move crash? \\
    C2 reward-only   & global NPC only    & on  & predicted to starve (3\% coverage) \\
    C3 full          & LLM director       & on  & the proposed method \\
    \bottomrule
  \end{tabular}
\end{table}

% TODO figures once runs land:
% \begin{figure}[H]\centering\includegraphics[width=0.85\textwidth]{imgs/arms_crash_bars}
%   \caption{Shield-off crash rate at 60 NPCs across the four arms.}\label{fig:cov-arms}\end{figure}

\section{Discussion}\label{sec:cov-discussion}

% Narrative to write:
%  - the experience > gradient > incentive hierarchy: TTC (incentive) failed for lack of
%    exposure; shield BC (gradient) fixed steering but not throttle; only when an LLM
%    manufactures the experience does a learned dense reward teach self-driven avoidance.
%  - salvage story: the abandoned meta-stack (schema-constrained decoding, deterministic
%    governor, async runner) is reused as trustworthy offline annotation infrastructure;
%    the negative result of Ch.7 is what FORCED this (correct) design, not a lucky guess.
%  - honest limitation: if the decisive test fails, report the bounded negative result.

\clearpage

%   line in main.tex AFTER the figures listed below exist in
%   Memoria/imgs/ and the bib keys in §10 of
%   docs/shielded_motif_design.md are added to referencias.bib.
%
%   Planned figures (imgs/, generate from run SQLite as in Ch.7):
%     - encounter_rate_curve   (exposure director lifts encounter_rate ~3% -> ~30%)
%     - reward_field           (r_phi over front-range x closing-speed slice)
%     - arms_crash_bars        (C0/C1/C2/C3 shield-OFF crash @60 NPC)
%     - shaping_learning_curve (shield-OFF crash vs updates, per arm)
%   Planned tables:
%     - tab:cov-arms           (factorial arms x outcomes)
% ============================================================

\chapter{Learning Collision Avoidance from LLM-Derived Rewards}\label{ch:llmreward}

Chapter~\ref{ch:dependence} closed on an explicit open problem. Behavioural cloning of the
shield's steering removed shield \emph{dependence} \textemdash\ the weaned agent drives as well
unaided as it does shielded \textemdash\ but it left a residual \emph{collision} gap: without the
shield the agent still crashes on roughly one episode in six, because the shield's throttle is an
emergency brake and is therefore a poor longitudinal teacher (Section~\ref{sec:dep-collapse}).
That chapter argued the gap is ``better posed as a reward-shaping problem \dots\ that lets the
agent learn its \emph{own} gentle avoidance''. This chapter takes up that problem, and in doing so
revisits \textemdash\ and finally vindicates \textemdash\ the use of a Large Language Model in the
training loop, after the negative result of Section~\ref{sec:dep-llm}.

The chapter is organised around a single thesis: the earlier LLM failure and the failure of a
hand-designed time-to-collision (TTC) penalty share \emph{one} root cause, and that cause dictates
the design that finally works. Section~\ref{sec:cov-diagnosis} establishes the diagnosis;
Section~\ref{sec:cov-method} derives the method (an LLM-learned dense reward trained on an
LLM-generated exposure curriculum); Sections~\ref{sec:cov-exposure}--\ref{sec:cov-shaping} detail
its three components; Section~\ref{sec:cov-results} reports the experiments; and
Section~\ref{sec:cov-discussion} places the result in the \emph{experience\,>\,gradient\,>\,
incentive} hierarchy that organises the whole of Chapters~\ref{ch:dependence}--\ref{ch:llmreward}.

\section{Diagnosis: one root cause behind two negative results}\label{sec:cov-diagnosis}

Two attempts to use a language model or a designed reward to close the avoidance gap failed, and
they failed identically.

\paragraph{The permissiveness dial (Section~\ref{sec:dep-llm}).} The LLM controlled a single
scalar every hundred episodes; its output was \emph{not causally coupled} to the metric it was
meant to move (the correlation between permissiveness and shield-activation rate was
\(\approx 0\)). A fixed schedule would have done the same.

\paragraph{The TTC penalty.} A hand-designed time-to-collision / following-distance penalty was
added to the reward shaper and its weight swept across four runs. The penalty \emph{fired in only
\(\approx 3\%\) of episodes}, its magnitude when fired was flat across weights, and the shield-off
crash rate did not fall \textemdash\ it merely redistributed failures from off-road to collision.
% TODO: cite/forward-ref the TTC sweep table if it is written up elsewhere.

\paragraph{The common cause.} Both signals were starved. On the Town04 highway loop, globally
spawned traffic places a vehicle in the ego's front cone in only about \(3\%\) of episodes, so any
collision-specific signal \textemdash\ whether an LLM dial or a TTC penalty \textemdash\ almost
never has anything to act on. Three constraints follow, and the method of
Section~\ref{sec:cov-method} is the minimal design that satisfies all three:
\begin{enumerate}
  \item \textbf{Bandwidth.} The learning signal must be \emph{dense} (per step), like the
        behavioural-cloning term of Equation~\eqref{eq:bc}, not a per-window scalar.
  \item \textbf{Exposure.} A collision signal starves at \(3\%\) coverage; the relevant
        experience must be \emph{manufactured}.
  \item \textbf{Stationarity.} A live additive reward rescales the running return
        normaliser the critic is calibrated against (Section~\ref{sec:design-agent-ppo}); the
        signal must be \emph{policy-invariant}.
\end{enumerate}

\section{Method: dense LLM reward over generated exposure}\label{sec:cov-method}

The proposed method, which we call \emph{Shielded Motif under generated exposure}, has three
components that map one-to-one onto the three constraints. An \textbf{exposure director}
(Section~\ref{sec:cov-exposure}) manufactures vehicle encounters, supplying the missing
experience. A \textbf{preference-trained reward model} (Section~\ref{sec:cov-reward}) supplies a
dense, per-step safety signal learned from a language model's judgements \textemdash\ offline,
so the model never runs inside the control loop. A \textbf{potential-based shaping}
integration (Section~\ref{sec:cov-shaping}) delivers that signal without altering the optimal
policy. Throughout, the safety shield and the deterministic safety governor of
Chapter~\ref{ch:dependence} are retained unchanged: the shield still prevents unsafe
\emph{exploration} on the harder, manufactured scenarios, and the ``proposer~+~deterministic
veto'' pattern \textemdash\ already shown to be sound independently of the LLM
(Section~\ref{sec:dep-discussion}) \textemdash\ now guards the exposure curriculum.

The design follows the Motif paradigm of intrinsic motivation from artificial-intelligence
feedback~\cite{klissarov2023motif}, in which an LLM expresses preferences over pairs of
observation captions that are distilled into a reward model; it inherits the preference-learning
formulation of~\cite{christiano2017preferences} and the policy-invariance guarantee
of~\cite{ng1999shaping}. It is related to, but distinct from, language-conditioned reward
generation for driving~\cite{huang2025vlmrl} and to LLM-based safety-critical scenario
generation~\cite{zhang2024chatscene,nie2025seeking}, from which the exposure director borrows its
intent-conditioned design.

\section{Component I \textemdash\ the exposure director}\label{sec:cov-exposure}

% Narrative to write:
%  - route-traffic injection along the ego waypoint chain (leads + cut-ins),
%    absolute desired speeds (not relative to the speed limit), governor-gated.
%  - failure-driven scenario generation: shield-off crash traces -> LLM infers a
%    behavioural intent -> structured scenario spec; intent->recipe memory bank
%    (after nie2025seeking, zhang2024chatscene).
%  - progress KPI: encounter_rate = fraction of episodes with min_front_dynamic < 0.7.
%  - safety: a dedicated governor with exposure ceilings (crash 0.20 / target 0.12)
%    that CUTS traffic on a crash breach.
%  - figure: encounter_rate_curve.

\section{Component II \textemdash\ the preference-trained reward model}\label{sec:cov-reward}

% Narrative to write:
%  - offline pipeline: sample observations (focus on the dynamic-LIDAR channel
%    obs[240:480] + front-cone scalars), template each into a structured caption,
%    LLM ranks PAIRS into preferences (reuse the schema-constrained, deterministic,
%    fallback-guarded decoding of the abandoned shield tuner).
%  - reward model r_phi: MLP 240(+k)->64->64->1; Bradley-Terry loss below.
%  - DISTINCTION from a crash-label safety critic: dense graded preference signal
%    available on near-miss steps that never crashed, vs sparse binary crash labels.

The reward model \(r_\phi\) is trained from preferences with the Bradley--Terry
objective~\cite{christiano2017preferences}: for an annotated pair of scenes \((s_a, s_b)\) in
which the language model judges \(s_a\) the safer situation,
\begin{equation}\label{eq:bt}
  \mathcal{L}(\phi) \;=\; -\,\mathbb{E}_{(s_a,s_b)}\;
  c_{ab}\,\log \sigma\!\bigl(r_\phi(s_a) - r_\phi(s_b)\bigr),
\end{equation}
where \(\sigma\) is the logistic function and \(c_{ab}\in[0,1]\) is the language model's stated
confidence, used to down-weight uncertain pairs. The annotation is performed entirely offline, so
the language model never contends with CARLA for GPU memory and adds no latency to the \(20\,
\mathrm{Hz}\) control loop \textemdash\ directly removing the operational difficulty that the
in-loop controller of Section~\ref{sec:dep-llm} suffered.

\section{Component III \textemdash\ policy-invariant shaping}\label{sec:cov-shaping}

The learned reward is delivered as potential-based shaping~\cite{ng1999shaping}, using \(r_\phi\)
as the potential:
\begin{equation}\label{eq:potential}
  F_t \;=\; \gamma\,r_\phi(s_{t+1}) - r_\phi(s_t),
  \qquad
  r^{\mathrm{shaped}}_t \;=\; r_t \;+\; \beta\,F_t .
\end{equation}
Because \(F_t\) is a difference of potentials, the set of optimal policies is provably
unchanged~\cite{ng1999shaping}: the term can only \emph{re-time} credit toward the states that
precede danger, it cannot create a new optimum. This is the formal guarantee that the
park-attractor of Section~\ref{sec:dep-collapse} \textemdash\ in which a direct supervised
gradient overpowered the reward and drove the policy to a stationary equilibrium \textemdash\
cannot recur here, and it is the reason the design respects the reward-stationarity constraint
that the original meta-controller was deliberately built to honour. The coefficient \(\beta\) is
frozen per run and annealed toward a small floor, strong early to install the avoidance skill.

\section{Experiments}\label{sec:cov-results}

% Pre-registered factorial design (see docs/shielded_motif_design.md, table tab:cov-arms):
%  C0 control (= learn_fix2), C1 exposure-only, C2 reward-only (predicted to STARVE),
%  C3 full. All warm-start learn_fix2 with --log_std_anchor_coef 1e-3 (mandatory).
%  Primary metric: shield-OFF crash @60 NPC (main_eval --num_npc 60 --seed 100).
%  Success: C3 shield-OFF crash < 16% (the learn_fix2 control).
%  Decisive/falsification test: if exposure lifts encounter_rate to ~30% and crash
%  still will not move, that is a clean, causally-isolated third negative result.

\begin{table}[H]
  \centering
  \caption{Pre-registered factorial arms. Each isolates one leg of the method; C2 is predicted
           to \emph{starve} (a reward with no exposure), which would itself confirm the
           exposure diagnosis of Section~\ref{sec:cov-diagnosis}.}
  \label{tab:cov-arms}
  \small
  \begin{tabular}{@{}llll@{}}
    \toprule
    \textbf{Arm} & \textbf{Exposure} & \textbf{Pref.\ reward} & \textbf{Purpose} \\
    \midrule
    C0 control       & global NPC only    & off & reproduce \texttt{learn\_fix2} \\
    C1 exposure-only & LLM director       & off & does exposure alone move crash? \\
    C2 reward-only   & global NPC only    & on  & predicted to starve (3\% coverage) \\
    C3 full          & LLM director       & on  & the proposed method \\
    \bottomrule
  \end{tabular}
\end{table}

% TODO figures once runs land:
% \begin{figure}[H]\centering\includegraphics[width=0.85\textwidth]{imgs/arms_crash_bars}
%   \caption{Shield-off crash rate at 60 NPCs across the four arms.}\label{fig:cov-arms}\end{figure}

\section{Discussion}\label{sec:cov-discussion}

% Narrative to write:
%  - the experience > gradient > incentive hierarchy: TTC (incentive) failed for lack of
%    exposure; shield BC (gradient) fixed steering but not throttle; only when an LLM
%    manufactures the experience does a learned dense reward teach self-driven avoidance.
%  - salvage story: the abandoned meta-stack (schema-constrained decoding, deterministic
%    governor, async runner) is reused as trustworthy offline annotation infrastructure;
%    the negative result of Ch.7 is what FORCED this (correct) design, not a lucky guess.
%  - honest limitation: if the decisive test fails, report the bounded negative result.

\clearpage
